Es llança un dau. Si surt 1 o 2, es treu una bola de l'urna $U_1$, que té 3 boles blanques i 2 de negres; si
no, es treu de l'urna $U_2$, que té 1 bola blanca i 4 de negres.

\begin{apartats}

\begin{tria}{bayes-tasca}
\itemtria{original}{0,75}{1,25}
Quina és la probabilitat que la bola sigui blanca?

\begin{solucio}
$P(U_1)=\dfrac26=\dfrac13$ i $P(U_2)=\dfrac23$. Per la probabilitat total,
$P(B)=\dfrac13\cdot\dfrac35+\dfrac23\cdot\dfrac15=\dfrac3{15}+\dfrac2{15}=\dfrac13$.
\end{solucio}

\itemtria{composicio-urna}{0,75}{1,25}
Quantes boles blanques hauria de tenir l'urna $U_2$, si continués tenint 5 boles en total, perquè la
probabilitat de treure una bola blanca fos $\dfrac7{15}$?

\begin{solucio}
Si $U_2$ tingués $b$ boles blanques, seria
$P(B)=\dfrac13\cdot\dfrac35+\dfrac23\cdot\dfrac b5=\dfrac15+\dfrac{2b}{15}$.\\
Cal $\dfrac15+\dfrac{2b}{15}=\dfrac7{15}$, és a dir $3+2b=7$: $\boxed{b=2}$ boles blanques (i 3 de
negres).
\end{solucio}
\end{tria}

\apartat[1,25]{1}
Si la bola és blanca, quina és la probabilitat que sigui de $U_1$? I si és negra, quina és la probabilitat
que sigui de $U_2$?

\begin{solucio}
$P(U_1\mid B)=\dfrac{\frac13\cdot\frac35}{\frac13}=\dfrac35$.\\
$P(N)=1-\dfrac13=\dfrac23$, i $P(U_2\mid N)=\dfrac{\frac23\cdot\frac45}{\frac23}=\dfrac45$.
\end{solucio}

\begin{nomesllarg}
\apartat{0,75}
Ara, després de llançar el dau, es treuen dues boles de l'urna que toqui, sense reemplaçament. Quina és la
probabilitat que totes dues siguin blanques?

\begin{solucio}
De $U_1$: $\dfrac35\cdot\dfrac24=\dfrac3{10}$. De $U_2$ és impossible, perquè només hi ha una bola blanca.
Per tant, $P=\dfrac13\cdot\dfrac3{10}+\dfrac23\cdot0=\dfrac1{10}$.
\end{solucio}
\end{nomesllarg}

\end{apartats}
